% !TEX program = pdflatex
\documentclass[11pt,a4paper]{article}
\usepackage[margin=24mm]{geometry}
\usepackage{booktabs,tabularx,array}
\usepackage{amsmath,amssymb}
\usepackage{hyperref}
\usepackage{enumitem}
\usepackage{url}
\hypersetup{colorlinks=true,linkcolor=blue,urlcolor=blue}
\title{研究計画書\\
\large 参照条件付き生成映像制作における品質・工程コストの実証研究}
\author{Funding proposal draft --- bonsai/r2v}
\date{2026年10月18日}
\begin{document}
\maketitle

\section*{1. 研究概要}
本研究は、参照画像に基づく動画生成（Reference-to-Video; R2V）をアニメーション制作工程に組み込む際、どの制御・評価要素が品質改善と費用対効果に寄与するかを実験的に明らかにする。特に、棒人間／ポーズ制御と、命題・証拠・評価基準を用いるJEV型評価・修正ループを小規模な比較実験で検証する。生成モデルの新規開発を初期目標とはせず、既存技術を制作工程に適用し、不要な複雑さや費用を除外する方法を研究する。

\section*{2. 背景と必要性}
R2V研究は、参照対象の同一性とテキスト指示への適合を同時に扱い、動画生成評価研究は品質の多次元評価と人間判断との整合を検討している。しかし、モデルの評価スコアが、実制作での採用可否、手戻り、レビュー時間、採用ショットあたりの総費用にそのまま対応するとは限らない。制作工程においては、出力の品質に加え、失敗を早期に発見し、原因を特定し、過剰な再生成を防ぐ仕組みが必要となる可能性がある。

TPSのJidoka（異常を検出して不良の流出を防ぐ）やJust-in-Time（必要なものを必要な時に流す）は、品質ゲートや待ち時間削減を考える参照枠となる。ただし、創造的制作にTPSをそのまま移植できるとは仮定しない。短い評価・修正サイクルなどの代替方式と比較し、品質と総費用の実測に基づいて採用可否を判断する。

\section*{3. 企業・現場の課題仮説}
本計画では以下を検証対象となる仮説として扱い、企業全般の確定した事実とはみなさない。
\begin{itemize}[nosep]
\item 参照画像があっても、キャラクターの外見・衣装・対象の関係がショットごとに変化する。
\item 指定動作や構図が再現されず、再生成や人手修正が増える。
\item 自動評価スコアと制作担当者の採用判断が一致せず、見た目の破綻を合格としてしまう危険がある。
\item 生成条件や採否理由の記録不足により、成功例の再現と失敗原因の追跡が難しくなる。
\item 品質ゲートの導入は不良の流出を防ぐ一方、レビュー待ちや工程負荷を増やす可能性もある。
\end{itemize}
課題の妥当性は、公開資料の確認と、可能であれば制作実務者への半構造化インタビューによって検討する。インタビューが実施できない場合は、その制約を明記し、企業一般への主張を行わない。

\section*{4. 研究目的と研究質問}
目的は、品質の最低条件を満たす構成の中から、最も費用対効果の高い工程条件を特定することである。
\begin{description}[style=nextline,leftmargin=0pt]
\item[RQ1] 棒人間／ポーズ制御は、同じR2Vベースラインに比べて指定動作への一致を改善し、追加費用に見合うか。
\item[RQ2] ポーズ制御にJEV型評価・修正ループを加えると、初回合格率、手戻り、人間レビュー時間は改善するか。
\item[RQ3] 人間の最低品質条件を満たした上で、どの条件が採用ショットあたり総費用とリードタイムの面で最も効率的か。
\item[RQ4] JEVと人間の判断はどの項目で一致・不一致となり、誤PASSをどの程度抑えられるか。
\end{description}

\section*{5. 仮説}
\begin{itemize}[nosep]
\item[H1.] 棒人間／ポーズ制御は、通常生成に比べて指定動作への一致を改善する可能性があるが、すべてのショットで追加費用に見合うとは限らない。
\item[H2.] JEV型の評価・修正ループは、明示的に検証できる項目では手戻りを減らす可能性がある。一方、人間の美的・演技的判断との不一致が残る可能性がある。
\item[H3.] 人間レビューを含む総費用で比較すると、すべての構成要素を導入する方式より、必要な要素だけを用いる構成が有利なタスク群が存在する可能性がある。
\item[H4.] 背景の物体数・位置・色などの明示的な整合性は自動判定候補となり得るが、顔、手、演技、感情、画風、構図意図は人間評価を残す方が安全である。
\end{itemize}

\section*{6. 研究方法}
\subsection*{6.1 初期実験の3条件}
\begin{center}
\begin{tabularx}{\linewidth}{cX}
\toprule
条件 & 内容 \\
\midrule
A & R2Vベースラインによる通常生成 \\
B & Aに棒人間／ポーズ制御を追加 \\
C & BにJEV型評価・修正ループを追加 \\
\bottomrule
\end{tabularx}
\end{center}
YAMLは仕様、モデル版、参照素材、生成条件、評価結果を記録する共通基盤として全条件で使用し、初期実験ではYAML自体の独立効果を主張しない。LoRA、追加MLアダプター、TPS型品質ゲート、Kaggleでの大規模実験、Kubernetesは初期条件に含めない。初期結果または失敗分析が必要性を示した場合に限り、後続の独立実験として追加する。

\subsection*{6.2 課題ショットとデータ管理}
小規模な課題セットを作成し、単純な動作と難しい動作、単一人物と複数対象、背景整合性、参照対象の一貫性を含める。課題数・試行数は利用可能な計算資源と事前のパイロット試験を基に決め、結果を見た後に都合よく変更しない。参照画像・動画の出典、ライセンス、利用目的、モデル規約を記録する。学習・調整に用いる素材と最終評価素材は分離する。

\subsection*{6.3 評価と盲検化}
評価ルーブリックを実験前に固定する。人間評価者には条件名を伏せ、同じ尺度で出力を評価してもらう。可能な範囲で複数評価者を用い、評価者間一致度を報告する。JEVは、命題、証拠、判定基準、確信度、失敗理由を記録する。顔・手・演技・感情・画風・構図意図については、JEVの合格だけで最終採用を決めない。

\section*{7. 評価指標}
\begin{itemize}[nosep]
\item \textbf{品質}: 指定動作への一致、参照対象の一貫性、時間的連続性、指示適合、人間評価による最低品質条件。
\item \textbf{JEV}: 人間評価との一致度、誤PASS率（人間不合格をJEVが合格とした割合）、誤REJECT率。
\item \textbf{工程}: 初回合格率、再生成回数、手戻り率、レビュー時間、リードタイム。
\item \textbf{費用}: 推論費、学習・調整費、評価器実行費、再生成費、人間レビュー・後処理工数。
\end{itemize}
主たる費用指標は次とする。
\[
\mathrm{CostPerAcceptedShot} =
\frac{\mathrm{Generation}+\mathrm{Training}+\mathrm{Evaluation}+\mathrm{Review}+\mathrm{Rework}}
{\mathrm{HumanAcceptedShots}}
\]
品質最低条件を満たさない構成は、費用が低くても最適とは扱わない。効果量と不確実性を報告し、統計的有意差の有無だけで採否を決めない。

\section*{8. 分析計画}
A/B/C間で同じ課題、同じ評価基準、可能な範囲で同じモデル版と生成設定を用いる。課題ごとの対応を保ち、平均値だけでなく分布、信頼区間、失敗モードを報告する。評価者が少数の場合は、統計的な一般化に限界があることを明示する。改善が見られない場合も重要な結果として扱い、試行数不足と実用上の効果不足を区別する。

\section*{9. 期待される成果}
\begin{enumerate}[nosep]
\item R2V制作を対象にした、再現可能な小規模比較プロトコル。
\item 棒人間／ポーズ制御とJEV型修正ループの品質・費用への限界寄与の測定。
\item 自動評価と人間評価の不一致を分類する評価データ。
\item どの要素を導入し、どの要素を外すべきかを判断する費用対効果の基準。
\item サーベイ論文、研究計画書、実験ログ、評価ルーブリック、再現用コードの公開可能な成果物。
\end{enumerate}
成果は仮説であり、改善や費用削減を事前に保証しない。

\section*{10. 実施スケジュール（6か月案）}
\begin{center}
\begin{tabularx}{\linewidth}{p{18mm}X}
\toprule
期間 & 作業と成果物 \\
\midrule
1か月目 & 文献調査の検索式・採否基準の固定、企業課題仮説の確認、ライセンス確認 \\
2か月目 & 課題ショット選定、評価ルーブリック、YAMLスキーマ、実験計画の事前固定 \\
3か月目 & ベースライン実装、パイロット、実行コスト計測の検証 \\
4か月目 & A/B/C比較実験、生成ログ・失敗例の収集 \\
5か月目 & 盲検評価、JEVとの照合、費用・不確実性分析 \\
6か月目 & 再現確認、結果整理、論文・報告書・公開可能なデータの整備 \\
\bottomrule
\end{tabularx}
\end{center}

\section*{11. 予算案（概算・申請先の規程に合わせて要調整）}
以下は予算枠を相談するための仮置きであり、見積書・謝金規程・助成条件を確認した正式な見積もりではない。研究者本人の人件費を計上できるか、クラウド利用料・謝金の扱い、間接経費率は申請先ごとに確認する。
\begin{center}
\begin{tabularx}{\linewidth}{p{37mm}p{20mm}X}
\toprule
費目 & 仮置き額 & 積算の考え方 \\
\midrule
クラウドGPU・実験実行 & 30万円 & 小規模推論、比較実験、再実行の上限枠 \\
評価協力者・謝金 & 25万円 & 盲検評価、評価者間一致の確認（人数・時間はパイロット後に確定） \\
研究補助・データ整理 & 30万円 & 課題素材整理、ログ確認、アノテーション補助 \\
データ保管・ソフトウェア等 & 5万円 & 必要な有料サービス、保管、実験管理の実費枠 \\
成果公開・予備費 & 10万円 & 公開・再現確認・予期せぬ再実行に備える枠 \\
\midrule
直接経費合計（案） & 100万円 & 申請先の対象費目と相見積もりに応じて調整 \\
\bottomrule
\end{tabularx}
\end{center}
GPU利用は無料枠・公開モデル・手元での軽量試験を優先し、有料計算は必要性が確認できた場合のみ利用する。追加学習や大規模データ生成を初期段階から予算化しない。間接経費、消費税、旅費、研究者人件費は申請制度に応じて別途積算する。

\section*{12. リスク管理と代替策}
\begin{itemize}[nosep]
\item \textbf{計算資源不足}: 軽量モデル、短いクリップ、少数ショットでベースラインを成立させる。Kaggle等は候補環境であり必須条件ではない。
\item \textbf{評価者の確保が難しい}: まず少人数のパイロットでルーブリックを修正し、結論の一般化範囲を制限する。
\item \textbf{参照素材の権利問題}: ライセンス確認済みの公開素材、研究利用が許可された素材、独自作成素材を優先する。
\item \textbf{JEVが人間判断に一致しない}: 失敗を隠さず、誤PASSの発生条件と自動化に向かない項目を報告する。
\item \textbf{TPS型ゲートが遅延を生む}: 方式を固定せず、短いフィードバックループ等の代替案を比較対象にする。
\item \textbf{改善効果が小さい}: 否定的結果も成果として、不要な要素を除外するための証拠と限界を報告する。
\end{itemize}

\section*{13. 倫理・再現性}
著作権・利用規約・肖像・個人情報に配慮し、利用可能なデータだけを使用する。人間評価では目的、利用範囲、保存期間を説明し、必要な同意を得る。個人を識別する素材を扱う場合は匿名化・アクセス制限を検討する。モデル、データ、プロンプト、設定、評価基準、実行ログをバージョン管理し、再現不能な箇所は明示する。

\section*{14. 研究上の位置づけ}
本研究は、R2Vモデル自体の新規開発よりも、既存生成技術を制作工程に組み込む際の評価・制御・費用測定に焦点を当てる。新規性は現段階で確定したものではなく、先行研究の体系的調査を通じて差分を確認する。TPS、Software Factory、JEV、棒人間制御を一括して導入するのではなく、初期実験を3条件に限定することで、何が寄与し何が不要かを検証する。

\end{document}
